% Movimento Dependente: 3 blocos e 4 roldanas 
% Faro, dom 07 dez 2025 15:37:47 
% Written by Tordar 
% pstricks2pdf.sh 3blocos_4roldanas
\documentclass{article}
\pagestyle{empty}

\usepackage{pst-all} 
\usepackage{pst-slpe}
\usepackage{graphicx}
\input{apnum}

% Roldana
\newcommand{\roldana}[5]{% x, y, radius, hole-radius, color
 \pscircle[linewidth=0.5mm,linecolor=lightgray,dimen=middle,fillstyle=radslope,slopebegin=gray,slopeend=white,slopesteps=50](#1,#2){#3}
  \ifnum\pdfstrcmp{#4}{0}=0\relax\else
    \pscircle[fillcolor=white,fillstyle=solid,linewidth=0.8pt](#1,#2){#4}
  \fi
  \pscircle[linewidth=1pt,linecolor=#5!50!black](#1,#2){#3}
}
% Bloco
\newcommand{\bloco}[5][]{% #1=options, #2=cx, #3=cy, #4=totalWidth, #5=height
  % Left frame (half of total width, full height)
  \psframe[linewidth=1pt,linecolor=gray,fillstyle=slope,slopebegin=white,slopeend=gray,slopesteps=15]
    (!#2 #4 2 div sub #3 #5 2 div sub)                     % lower-left of total area
    (!#2 #3 #5 2 div add)                                  % upper-right of left frame (center x, top)  
  % Right frame (half of total width, full height)
  \psframe[linewidth=1pt,linecolor=gray,fillstyle=slope,slopebegin=gray,slopeend=white,slopesteps=15]
    (!#2 #3 #5 2 div sub)                                  % lower-left of right frame (center x, bottom)
    (!#2 #4 2 div add #3 #5 2 div add)                     % upper-right of total area
  \psframe[#1](!#2 #4 2 div sub #3 #5 2 div sub)(!#2 #4 2 div add #3 #5 2 div add) 
}

\begin{document}

% Raio de cada roldana:
\evaldef\ther{0.75}
% Altura e largura de cada bloco: 
\evaldef\thew{1.5}
\evaldef\theh{1.5}
% Coordenadas da corda: 
\evaldef\xi{2}
\evaldef\yi{4}
\evaldef\xii{\xi+\ther}
\evaldef\yii{\yi-\ther}
\evaldef\xiii{\xii+\ther}
\evaldef\yiii{\yi}
\evaldef\xiv{\xiii}
\evaldef\yiv{7}
\evaldef\xv{\xiv+\ther}
\evaldef\yv{\yiv+\ther}
\evaldef\xvi{\xv+\ther}
\evaldef\yvi{\yv-\ther}
\evaldef\xvii{\xvi}
\evaldef\yvii{3}
\evaldef\xviii{\xvii+\ther}
\evaldef\yviii{\yvii-\ther}
\evaldef\xix{\xviii+\ther}
\evaldef\yix{\yviii+\ther}
\evaldef\xx{\xix}
\evaldef\yx{6}
\evaldef\xxi{\xx+\ther}
\evaldef\yxi{\yx+\ther}
\evaldef\xxii{\xxi+\ther}
\evaldef\yxii{\yxi-\ther}
\evaldef\xxiii{\xxii}
\evaldef\yxiii{4}
% Coordenadas da primeira roldana: 
\evaldef\xri{\xii}
\evaldef\yri{\yi}
% Coordenadas da segunda roldana: 
\evaldef\xrii{\xv}
\evaldef\yrii{\yiv}
% Coordenadas da terceira roldana: 
\evaldef\xriii{\xviii}
\evaldef\yriii{\yvii}
% Coordenadas da quarta roldana: 
\evaldef\xriv{\xxi}
\evaldef\yriv{\yx}
% Coordenadas do primeiro bloco: 
\evaldef\xbi{\xri}
\evaldef\ybi{2}
% Coordenadas do segundo bloco: 
\evaldef\xbii{\xriii}
\evaldef\ybii{1}
% Coordenadas do terceiro bloco: 
\evaldef\xbiii{\xriv+\ther}
\evaldef\ybiii{3.5}
% Coordenadas para fios roldana - teto:
\evaldef\xrti{\xrii}
\evaldef\yrti{\yrii}
\evaldef\xrtii{\xriv}
\evaldef\yrtii{\yriv}
% Coordenadas para fios roldana - bloco:
\evaldef\xrbi{\xri}
\evaldef\yrbi{\yri}\evaldef\yrbie{\ybi+\theh/2}
\evaldef\xrbii{\xriii}
\evaldef\yrbii{\yriii}\evaldef\yrbiie{\ybii+\theh/2}
% Coordenadas para identificação dos blocos: 
\evaldef\xba{\xbi}
\evaldef\yba{\ybi-0.75*\theh}
\evaldef\xbb{\xbii}
\evaldef\ybb{\ybii-0.75*\theh}
\evaldef\xbc{\xbiii}
\evaldef\ybc{\ybiii-0.75*\theh}

% Let's go:
\begin{pspicture}(0,0)(10,10)
% Teto: 
\psframe[linewidth=0pt,linecolor=lightgray,dimen=inner,fillstyle=hlines*,fillcolor=lightgray,hatchangle=45](1.5,9)(8.5,10) 
% Corda: 
\psline[linewidth=2pt,linecolor=black,linearc=\ther](\xi,9)(\xi,\yi)(\xii,\yii)(\xiii,\yiii)(\xiv,\yiv)(\xv,\yv)(\xvi,\yvi)(\xvii,\yvii)(\xviii,\yviii)(\xix,\yix)(\xx,\yx)(\xxi,\yxi)(\xxii,\yxii)(\xxiii,\yxiii)
% Roldanas:
\roldana{\xri}{\yri}{\ther}{0.05}{gray} % 1 
\roldana{\xrii}{\yrii}{\ther}{0.05}{gray} % 2
\roldana{\xriii}{\yriii}{\ther}{0.05}{gray} % 3
\roldana{\xriv}{\yriv}{\ther}{0.05}{gray} % 3
% Fios roldana - teto: 
\psline[linewidth=2pt,linecolor=black](\xrti,\yrti)(\xrti,9)
\psline[linewidth=2pt,linecolor=black](\xrtii,\yrtii)(\xrtii,9)
% Blocos: 
\bloco[linewidth=2pt,linecolor=black]{\xbi}{\ybi}{\thew}{\theh}
\bloco[linewidth=2pt,linecolor=black]{\xbii}{\ybii}{\thew}{\theh}
\bloco[linewidth=2pt,linecolor=black]{\xbiii}{\ybiii}{\thew}{\theh}
% Fios roldana - bloco:
\psline[linewidth=2pt,linecolor=black](\xrbi,\yrbi)(\xrbi,\yrbie) 
\psline[linewidth=2pt,linecolor=black](\xrbii,\yrbii)(\xrbii,\yrbiie)
\rput{0}(\xba,\yba){\scalebox{1.5}{$A$}}
\rput{0}(\xbb,\ybb){\scalebox{1.5}{$B$}}
\rput{0}(\xbc,\ybc){\scalebox{1.5}{$C$}}
\end{pspicture}

\end{document}
